\section{The 6-DOF path}\label{sec:sixdof}

The placement model defined in \cref{sec:datamodel} and resolved in \cref{sec:resolver} is, by
construction, a three-degree-of-freedom (3-DOF) system: every part is seated coaxially, child
origins differ from their parents only by an axial translation, and no relative orientation is
modeled. That restriction is not accidental --- it matches what the current integrator can use
(\cref{sec:intro}) --- but it is a deliberately \emph{soft} restriction. The data model was shaped
so that the eventual move to six degrees of freedom (6-DOF) costs \emph{no schema change, no
signature change, and no second resolver}. This section documents exactly which seams were widened
in advance, why each widening is free today, and what the one structural code change will be when
the rotational degrees of freedom are switched on. Nothing here is implemented or exercised in the
3-DOF cut; it is the forward-compatibility argument for the choices already made.

The motivating constraint is stated in the source itself. The composite mathematics today rests on
an explicit frame assumption: ``all parts share the same body-frame orientation, so child position
offsets are pure translations and tensors combine by addition (no rotation) \ldots\ relative part
rotation is not modeled'' (\src{model/parts/Part.h}{43}--45). The value of the 6-DOF shaping is that
it converts this from a structural assumption baked across the tree into a single localized
limitation --- one identity-valued quaternion and one composition line --- that can be lifted in
place.

\subsection{Orientation already lives in the value types}

Two carriers for orientation are present in the data model from day one, both holding the identity
quaternion in the 3-DOF cut so that behavior is bit-identical to a pure-translation tree.

The first is \code{Pose}, the resolver's output type (\cref{sec:resolver}). It pairs the part-local
origin with an orientation, and its \code{compose} operator already rotates the child translation by
the parent orientation rather than merely adding it:

\begin{listing}[htbp]
\begin{minted}{cpp}
struct Pose {
   Vector3    origin {Vector3::Zero()};        // aft-plane origin, in ROOT frame
   Quaternion orient {Quaternion::Identity()}; // (x,y,z,w); identity in 3-DOF

   Pose compose(const Pose& childInThis) const {
      return Pose{ origin + orient * childInThis.origin,
                   (orient * childInThis.orient).normalized() };
   }
};
\end{minted}
\caption{\code{Pose::compose} rotates the child translation by the parent orientation. With every
\code{orient} equal to identity, \code{orient * childInThis.origin} is the input unchanged, so
the composed pose is a pure translation --- bit-identical to the current
tree.}\label{lst:pose-compose}
\end{listing}

Because \code{orient} is \cppinline|Quaternion::Identity()| everywhere in 3-DOF, the rotation
\cppinline|orient * childInThis.origin| is the identity map and \code{compose} degenerates to vector
addition. The operator is written correctly for the rotated case now, so no edit to the resolver's
spine is required later; the rotations simply stop being identities.

The second carrier is \code{StationLink::childRot}, the per-seam orientation defined in the link
struct (\cref{sec:datamodel}). It is a \emph{value} in the already-existing struct, not a new field
introduced by the 6-DOF work --- the storage layout does not reshape. It defaults to
\cppinline|Quaternion::Identity()| and is consumed by the resolver as the child's local orientation
relative to its parent. This is where canted fins and side boosters --- parts that genuinely seat at
a fixed relative angle --- are expressed: a non-identity \code{childRot} on the seam, resolved
through \code{compose} into a world orientation, with no other part of the pipeline aware that
anything changed.

\begin{keyinsight}[Zero schema cost is the point]
Carrying \code{orient} on \code{Pose} and \code{childRot} on \code{StationLink} from the first
commit means the on-disk \code{.qrd} schema (\cref{sec:serialization}), the storage type for
\code{childParts}, and the resolver signature \cppinline|resolvePlacements(const Part&, const Pose&)|
are all already 6-DOF-shaped. Enabling rotation adds behavior, not types. A design that bolted
orientation on later would have to migrate every stored design and re-break every consumer of
\code{getChildParts()} a second time --- precisely the churn this whitepaper exists to avoid.
\end{keyinsight}

\subsection{The one new composition line: rotating the inertia tensor}

The single structural change that distinguishes 6-DOF from 3-DOF lives in the composite derivation,
not the resolver. The two-pass mass-weighted CM and parallel-axis routine in \code{computeCompositeAt}
(\src{model/parts/Part.cpp}{169}--204) is preserved verbatim by the placement refactor
(\cref{sec:resolver}); only its \emph{input} changes from stored CM-to-CM offsets to
resolver-derived geometric stations. The 6-DOF step adds exactly one operation to that preserved
routine: before a child's inertia tensor is shifted to the composite CM by the parallel-axis theorem
(\code{parallelAxisTerm}, \src{model/parts/Part.cpp}{17}--23), it must first be rotated out of the
child's body frame into the root frame.

\begin{listing}[htbp]
\begin{minted}{cpp}
// In computeCompositeAt's pass 2, per child, with R = childPose.orient:
const Matrix3 R    = childPose.orient.toRotationMatrix();
const Matrix3 Irot = R * childI * R.transpose();   // I' = R I R^T  (NEW for 6-DOF)
const Vector3 d    = childCmInRoot - compositeCm;   // child CM -> composite CM
I += Irot + childMass * parallelAxisTerm(d);        // parallel-axis shift, unchanged
\end{minted}
\caption{The lone 6-DOF addition to the composite math: the similarity transform
\(I' = R\,I\,R^{\mathsf T}\) re-expresses the child's centroidal tensor in the root frame before the
parallel-axis shift. In 3-DOF \(R = \mathbf{I}_3\), so \code{Irot == childI} and the line is an
identity.}\label{lst:tensor-rotate}
\end{listing}

The transform \(I' = R\,I\,R^{\mathsf T}\) is the standard congruence that re-expresses a centroidal
inertia tensor under a change of body orientation, where \(R\) is the rotation matrix of
\code{childPose.orient}. It must precede the parallel-axis shift because the displacement tensor
\code{parallelAxisTerm} operates on the offset \code{d} in the root frame; the child tensor must be
in that same frame before the two are summed. In the 3-DOF cut \(R\) is the identity, so \code{Irot}
equals \code{childI} bit-for-bit and the routine reduces to exactly the addition it performs today.
This is why \src{model/parts/Part.h}{43}--45 can be cited as the \emph{current} limit rather than a
permanent one: that ``no rotation'' assumption is now localized to this single line, and lifting it
is a one-line change to a routine whose structure is otherwise untouched.

\subsection{Off-axis radial seating arrives with rotation, not before}

\cref{sec:resolver} deliberately scopes the radial coordinate \(r\) to the seam and overlap
\emph{check} only: the pose solve sets \(x = y = 0\) and every part is seated coaxially. Off-axis
radial placement --- a single rail button, an asymmetric pod, a clustered side booster offset from
the centerline --- is not a 3-DOF feature held back arbitrarily. It is physically meaningful only
once two ingredients are present together: \code{childRot} (so the off-axis part is correctly
\emph{oriented} at its radial station, not merely translated to \(x = r\) while still pointing along
the parent axis) and a per-part body-frame center-of-mass offset.

The second ingredient is the latent purpose of \code{getCenterMassOffset()}
(\src{model/parts/Part.h}{96}). That accessor returns the part's \code{cm} member, which is
documented as not currently consumed by the composition math but kept as ``the natural home for that
offset once asymmetric parts or 6-DOF force application \ldots\ need it'' (\src{model/parts/Part.h}{310}--316).
In the 3-DOF cut its sole consumer is the cone CM-frame adapter \code{cmStationLocal}
(\cref{sec:datamodel}). Under 6-DOF it gains its second documented use: the radial components of the
body-frame CM offset, rotated by \code{childPose.orient} alongside the axial \code{cmStationLocal}
term, place a genuinely off-axis part's mass correctly in the root frame.

\begin{designrule}[Do not store a half-correct off-axis placement]
Seating a part at \(x = r\) in 3-DOF --- with no orientation and no body-frame CM offset --- would
record a placement that is correct in translation but wrong in attitude and inertia. That is exactly
the staleness-and-divergence error class this design exists to eliminate (\cref{sec:philosophy}).
Off-axis seating therefore lights up \emph{here}, with \code{childRot} and the body-frame CM offset
integrated together, and not one step sooner. Until then \(r\) drives only the diagnostic check
(\cref{sec:diagnostics}), where a value with no placement consequence is harmless.
\end{designrule}

\subsection{A body-frame integrator reuses the same resolver}

The final piece of the 6-DOF path requires no change to the placement subsystem at all. A future
body-frame integrator --- one that integrates the rocket's attitude as well as its translation ---
supplies the rocket's world orientation by multiplying the \code{rootPose} passed to
\cppinline|resolvePlacements(root, rootPose)| by the integrated attitude quaternion. Today
\code{rootPose} is conventionally the identity (the rocket's aft datum at the world origin,
\zfwd); a body-frame integrator passes the live attitude instead.

Because \code{compose} (\cref{lst:pose-compose}) already propagates that root orientation down the
entire tree, the very same resolver then yields every part's full world pose --- position and
orientation --- which the force and aero paths can consume to apply thrust and aerodynamic loads in
world coordinates. There is no schema change, no resolver-signature change, and no second code path:
the single resolver that serves the simulator and the visualizer in 3-DOF (\cref{sec:resolver})
serves the attitude-aware integrator unchanged. The 6-DOF work is thereby reduced to three localized
edits --- stop forcing \code{childRot} to identity at authoring and load time, rotate the child
tensor by the one line of \cref{lst:tensor-rotate}, and feed the integrated attitude into
\code{rootPose} --- against a data model and a resolver that were built to receive them.
